\section{Conclusion}

We have demonstrated the application of our broadband, multivariable PDE model on VUV sensitive devices. The PDE data for the Hamamatsu VUV4 and FBK VUV-HD technology was fit, yielding junction parameters, avalanche probabilities and oxide thicknesses allowing PDE prediction, validated against other reported measurements. Compared to the HPK-VUV4, the FBK VUV-HD's oxide is roughly 100 times thicker and the junction almost half the size. All measured $P_e$ values are similar, with a large spread in $P_h$ values. For operation in liquid nobles, we consider the PDE extrapolation at long wavelengths relevant to ExCT quite robust. Further model investigation and validation in the VUV, especially for LAr emission, is required. Immediate opportunities for expanding the model are measuring $P_e(V)$ and relative UV $PDE(\theta)$ over a wide temperature range. This model provides a foundation for improved detector simulation and SiPM characterization and understanding, and we hope combinations of the methods outlined here can be leveraged by other groups to characterize device PDE and extrapolate when measurement in the region of interest is not possible.  

The PDE model and lookup tables for simulation will be made available online.

%We anticipate future designs can push $FF$ close to unity while retaining high avalanche efficiency and timing resolution.

\section*{Acknowledgments}
The authors gratefully acknowledge support from Canadian Foundation for Innovation Fund (CFI) 2017. Additional support was provided by a grant from Canada Natural Sciences and Engineering Research Council of Canada (NSERC). We thank the Canadian taxpayers for supporting physics research.